%%%%%%%%%%%%%%%%%%%%%%%%%%%%%%%%%%%%%%%%%%%%%%%%%%%%%%%%%%%%%%%%%%%%%%
% Lớp 10 | Đại số tổ hợp | Nhị thức Newton
% Dạng 0D8?3-6: Toán tổ hợp có dùng nhị thức Newton
% Số câu: 11  (sinh tự động bởi tach_cau.py; từ nay thầy sửa trực tiếp file này)
%%%%%%%%%%%%%%%%%%%%%%%%%%%%%%%%%%%%%%%%%%%%%%%%%%%%%%%%%%%%%%%%%%%%%%

%%%%%%%%%%%%%%%% TRẮC NGHIỆM 4 PHƯƠNG ÁN %%%%%%%%%%%%%%%%

% nguồn: MaHoaNH_TracNghiem_Lop10
\begin{ex}%[0D8H3-6]%[ma:9fe1db25]
	 Tổng $T=\mathrm{C}_{2023}^0+\mathrm{C}_{2023}^1+\mathrm{C}_{2023}^2+\cdots+\mathrm{C}_{2023}^{2023}$ bằng
	 \choice
		{$2^{n+1}$} 
		 {$2^{2022}$}
		 {$2^n$}
		{\True  $2^{2023}$}
	 \loigiai{
		 Xét khai triển $(x+1)^n=\mathrm{C}_n^0+\mathrm{C}_n^1x+\mathrm{C}_n^2x^2+\cdots+\mathrm{C}_n^nx^n$.\\
		Với $x=1$ và $n=2023$,  ta có $T=\mathrm{C}_{2023}^0+\mathrm{C}_{2023}^1+\mathrm{C}_{2023}^2+\cdots+\mathrm{C}_{2023}^{2023}=2^{2023}$.
	 }
\end{ex}

% nguồn: MaHoaNH_TracNghiem_Lop10
\begin{ex}%[0D8V3-6]%[ma:bf01b022]
	Nếu một người gửi số tiền A vào ngân hàng theo thể thức lãi kép (đến kỳ hạn mà người gửi không rút lãi ra thì tiền lãi được tính vào vốn của kỳ kế tiếp) với lãi suất  r mỗi kì thì sau N kì, số tiền người ấy thu được cả vốn lẫn lãi là $C=A(1+r)^N$ (triệu đồng). Ông An
	gửi 20 triệu đồng vào ngân hàng $X$ theo thể thức lãi kép với lãi suất $8,65\%$ một quý. Hãy dùng ba số hạng đầu trong khai triển $(1+0,0865)^5$ tính sau 5 quý (vẫn tính lãi suất kì hạn theo quý),  ông An sẽ thu được số tiền cả vốn lẫn lãi là bao nhiêu (giả sử lãi suất hằng năm của ngân hàng X là không đổi)?
	 \choice
		{$30{,}15645$  triệu đồng}
		{\True $30{,}14645$ triệu  đồng}
		{$30{,}14675$ triệu đồng} 
		 {$31{,}14645$ triệu đồng}
	\loigiai{ 
		Áp dụng công thức $C=A(1+r)^5$ với $A=20$  triệu $r=8{,}65\%, n=5$ quý.
		$$ 
		\begin{aligned}
			& (1+x)^5=\mathrm{C}_5^0+\mathrm{C}_5^1  x+\mathrm{C}_5^2 x^2+\mathrm{C}_5^3 x^3+\mathrm{C}_5^4 x^4+\mathrm{C}_5^5 x^5. \\
			& (1+0{,}0865)^5 \approx \mathrm{C}_5^0+\mathrm{C}_5^1\cdot 0{,}0865+\mathrm{C}_5^2\cdot (0{,}0865)^2=1+5\cdot 0{,}0865+10  \cdot(0{,}0865)^2=1{,}5073225.
		\end{aligned}
		$$ 
		Vậy số tiền thu được sau $5$ quý là $C=20\cdot1{,}5073225=30{,}14645$ triệu đồng. 
	 }
\end{ex}

% nguồn: MaHoaNH_TracNghiem_Lop10
\begin{ex}%[0D8V3-6]%[ma:79898ba8]
	Để dự báo dân số của một quốc gia người ta sử dụng công thức $S=A(1+r)^n$, trong đó $A$ là dân số của năm  lấy làm mốc, $S$ là dân số sau $n$ năm, $r$ là tỉ lệ tăng dân số hàng năm, $r=1{,}5\%$. Năm $2015$ dân số của một quốc gia là $212\,942\,000$ người. Dùng ba số hạng đầu trong khai triển $(1+0{,}015)^5$ ta ước tính được số dân của quốc gia đó vào năm $2020$ gần số nào sau đây nhất?
	 \choice
		{\True $229\,391\,769$ nghìn người} 
		 {$329\,391\,769$ nghìn người}
		{$229\,391\,759$ nghìn  người}
		{$228\,391\,769$ nghìn  người}
	 \loigiai{
		Lấy năm $2015$ làm mốc và tính dân số năm $2015$  thì $n=2020-2015=5$.\\
		Áp dụng công thức $S=A(1+r)^n$ với  $A=212\,942\,000$, $r=1{,}5\%$.
		$$ 
		\begin{aligned}
			& (1+x)^5=\mathrm{C}_5^0+\mathrm{C}_5^1 x+\mathrm{C}_5^2 x^2+\mathrm{C}_5^3 x^3+\mathrm{C}_5^4 x^4+\mathrm{C}_5^5  x^5. \\
			& (1+0{,}015)^5 \approx \mathrm{C}_5^0+\mathrm{C}_5^1 \cdot 0{,}015+\mathrm{C}_5^2(0{,}015)^2=1+5 \cdot 0{,}015+10  \cdot(0{,}015)^2=1{,}07725.
		\end{aligned}
		 $$
		Ước tính dân số của quốc  gia đó vào năm $2020$ là $212\,942\,000\cdot 1{,}07725=229\,391\,769{,}5$.\\
		Vậy dân số quốc  gia đó là $229\,391\,769$ nghìn người.
	 }
\end{ex}


%%%%%%%%%%%%%%%% ĐÚNG/SAI %%%%%%%%%%%%%%%%

% nguồn: MaHoaNH_DungSai_Lop10
\begin{ex}%[0D8V3-6]%[ma:e6b45fa6]
	 Trong khai triển $ \left(x^2+2\right)^n $, biết $ A_n^3-8 C_n^2+C_n^1=49 $. Tìm các khẳng định đúng trong các khẳng định sau
	\choiceTF 
		{\True Giá trị của $n$ bằng  $7$}
		{\True Tổng hệ số trong khai triển là $280$} 
		{Hệ số của số hạng chứa $x^8$ trong khai triển là $560$} 
		{Hệ số của số hạng chứa $x^8$ trong khai triển là  $560x^8$}
	\loigiai{ 
		Điều kiện $ \heva{&n\ge  3\\&n\in\mathbb{N}} $.\\
		Phương trình đã cho tương đương với $\begin{aligned}
		 &\dfrac{n!}{(n-3)!}-8\cdot\dfrac{n!}{2!(n-2)!}+\dfrac{n!}{1!(n-1)!}=49\\
		\Leftrightarrow& (n-2)(n-1)n-4(n-1)n+n=49\\ 
		 \Leftrightarrow& n=7\,\text{(nhận)}\\
	\end{aligned}
	$\\ 
	Với  $ n=7 $ ta có $ \left(x^2+2\right)^n=\displaystyle\sum_{k=0}^{7}C_{7}^k\cdot(x^2)^{7-k}\cdot 2^k=\displaystyle\sum_{k=0}^{7}C_{7}^k\cdot 2^k\cdot x^{14-2k}$.\\
	Thay $x=1$ ta được tổng hệ  số là $(1+2)^7=280$.\\
	Để trong khai triển  đã cho có số hạng chứa $ x^8 $ là $ 14-2k=8\Leftrightarrow k=3$\,(nhận).\\
	Vậy số hạng của $ x^8 $ trong khai triển đã cho là $ \mathrm{C}_{7}^3\cdot 2^3\cdot x^8=280x^8$. 
 }
\end{ex}

% nguồn: MaHoaNH_DungSai_Lop10
\begin{ex}%[0D8V3-6]%[ma:190e7aa4]
	Trong khai triển $ \left(\dfrac{2}{x}-x^3\right)^n, \forall x \neq 0  $, biết $ \mathrm{C}_{n-4}^{n-6}+n \cdot \mathrm{A}_n^2=454 $.
	Tìm các khẳng định  đúng trong các khẳng định sau
	\choiceTF 
		{\True  Giá trị của $n$ bằng $8$}
		{\True Tổng hệ số trong khai triển  là $1$}
		{Hệ số của số hạng chứa $x^4$ trong khai triển  là $896$}
		{Hệ số của  số hạng chứa $x^4$ trong khai triển là $896x^4$}
	 \loigiai{
		Điều kiện $ \heva{&n\ge 6\\&n\in\mathbb{N}} $.\\ 
		Phương trình đã cho tương đương với $\begin{aligned}
		 &\dfrac{(n-4)!}{(n-6)!2!}+n\cdot\dfrac{n!}{(n-2)!}=454\\
		 \Leftrightarrow& \dfrac{(n-5)(n-4)}{2}+n\cdot(n-1)\cdot n=454\\
		 \Leftrightarrow&n=8\,\text{(nhận)}\\
	\end{aligned}
	 $\\
	Với $ n=8 $ ta có $ \left(\dfrac{2}{x}-x^3\right)^n =\displaystyle\sum_{k=0}^{8}C_{8}^k\cdot2^{8-k}\cdot x^{k-8}\cdot (x^3)^k=\displaystyle\sum_{k=0}^{8}C_{8}^k\cdot  2^{8-k}\cdot x^{4k-8}$.\\
	Thay $x=1$ ta được tổng  hệ số là $(2-1)^8=1$.\\
	Để trong khai triển đã cho có số hạng chứa $ x^4 $ là  $ 4k-8=4\Leftrightarrow k=3$\,(nhận).\\
	Vậy hệ số của $ x^4 $ trong khai triển đã  cho là $\mathrm{C}_{8}^3\cdot 2^5=1792$.
 }
\end{ex}


%%%%%%%%%%%%%%%% TRẢ LỜI NGẮN %%%%%%%%%%%%%%%%

% nguồn: MaHoaNH_TraLoiNgan_Lop10
\begin{bt}%[0D8H3-6]%[ma:9688dadf]
	Trền quầy có $5$ tờ vé số khác nhau. Một khách hàng có bao nhiêu lựa chọn  mua một số vé trong các vé số đó (tính cả không mua vé nào).
	 \shortans{32}
	 \loigiai{
	Số cách mua vé số là số tập con  của tập chứa $5$ phần tử nên số cách mua là
	$$ 
	\mathrm{C}^{5}_{5} + \mathrm{C}^{4}_{5} + \mathrm{C}^{3}_{5} + \mathrm{C}^{2}_{5} + \mathrm{C}^{1}_{5} +  \mathrm{C}^{0}_{5} = (1+1)^5 = 32.
	$$ 
	 }
\end{bt}

% nguồn: MaHoaNH_TraLoiNgan_Lop10
\begin{bt}%[0D8H3-6]%[ma:680104a7]
	Lớp $10$A đề nghị các tổ chọn thành viên để tập kịch. Tổ $I$ phải chọn ít nhất một thành viên để tham gia đội kịch của lớp.  Biết tổ $I$ có $5$ bạn, tìm số cách chọn thành viên để tập kịch của tổ $I$.
	\shortans{31} 
	\loigiai{ 
	Vì tổ $I$ cần chọn ít nhất một thành viên từ $5$ thành viên tổ nên số cách chọn thành  viên của tổ $I$ là
	 $$
	\mathrm{C}^{5}_{5} + \mathrm{C}^{4}_{5} +  \mathrm{C}^{3}_{5} + \mathrm{C}^{2}_{5} + \mathrm{C}^{1}_{5}
	= (1+1)^5 - \mathrm{C}^{0}_{5}  = 31.
	 $$
	 }
\end{bt}

% nguồn: MaHoaNH_TraLoiNgan_Lop10
\begin{bt}%[0D8C3-6]%[ma:0d8e93b9]
	Cho tập hợp $S$ có $n$  ($n\ge 3 $) phần tử. Gọi $A_1$, $A_2$, $A_3$, $A_4$, $A_5$, $A_6$ là các tập con của tập hợp $S$ thỏa mãn đồng thời
	\begin{itemize}
		\item $\left|A_i\right|=5$, với mọi $1\le i\le  6 $,
		\item $  \left|A_i\cap A_j\right| \le 2$, với mọi $1\le i<j\le 6 $.
	\end{itemize}
	Tìm giá trị nhỏ nhất của  $ n $.
	\shortans{10} 
	\loigiai{ 
		Đặt $ S=\left\{x_1;x_2;\ldots;x_n\right\} $. Ta đếm số bộ $\left(\left\{A_i, A_j\right\}, x_k\right)$, trong đó $x_k \in A_i \cap A_j$ và  $\left\{A_i, A_j\right\}$ không có tính thứ tự. Gọi số bộ này là $T$.\\
		Đếm theo  cặp tập hợp $\left(A_i, A_j\right)$, ta có $T \le 2 \mathrm{C}_6^2=30$.\\
		Đếm phần tử $x$. Gọi  $a_i$ là số tập hợp chứa phần tử $x_i$ với mọi $i=\overline{1, n}$. Khi đó, đếm theo số phần tử, ta có $a_1+a_2+\ldots+a_n=6\cdot 5=30$. Ngoài ra,
		$$ 
		T=\displaystyle \sum_{i=1}^n \mathrm{C}_{a_1}^2=\dfrac{1}{2}\left(a_1^2+a_2^2+\ldots+a_n^2-30\right) \ge  \dfrac{1}{2}\left[\dfrac{\left(a_1+a_2+\ldots+a_n\right)^2}{n}-30\right]=\dfrac{1}{2}\left(\dfrac{900}{n}-30\right).$$
		Suy ra $\dfrac{1}{2}\left(\dfrac{900}{n}-30\right)  \le 30 \Rightarrow n \ge 10$.\\
		Với $n=10$, ta có bảng sau  (ô được đánh dấu $x$ nghĩa là phần tử $x_{i} \in A_{j}$):
		\begin{center}
			\begin{tabular}{|c|c|c|c|c|c|c|c|c|c|c|}
				\hline & $\mathrm{x}_1$ & $\mathrm{x}_2$ & $\mathrm{x}_3$ & $\mathrm{x}_4$ & $\mathrm{x}_5$ & $\mathrm{x}_6$  & $\mathrm{x}_7$ & $\mathrm{x}_8$ & $\mathrm{x}_9$ & $\mathrm{x}_{10}$ \\
				\hline $A_1$ &  $\mathrm{x}$ & $\mathrm{x}$ & $\mathrm{x}$ & $\mathrm{x}$ & $\mathrm{x}$ & & & & & \\
				\hline $A_2$ & $\mathrm{x}$ & $\mathrm{x}$ & & & & $\mathrm{x}$ & $\mathrm{x}$ & $\mathrm{x}$ & &  \\
				\hline $A_3$ & $\mathrm{x}$ & & $\mathrm{x}$ & & & $\mathrm{x}$  & & & $\mathrm{x}$ & $\mathrm{x}$ \\
				\hline $A_4$ & & $\mathrm{x}$ & &  $\mathrm{x}$ & & & $\mathrm{x}$ & & $\mathrm{x}$ & $\mathrm{x}$ \\
				\hline $A_5$ & & & & $\mathrm{x}$ & $\mathrm{x}$ & $\mathrm{x}$ & & $\mathrm{x}$  & $\mathrm{x}$ & \\
				\hline $A_6$ & & & $\mathrm{x}$ & & $\mathrm{x}$ & & $\mathrm{x}$ & $\mathrm{x}$  & & $\mathrm{x}$ \\
				\hline 
			\end{tabular}
		\end{center}
		Vậy giá trị nhỏ nhất của $n$  là $10$.
	 }
\end{bt}


%%%%%%%%%%%%%%%% TỰ LUẬN %%%%%%%%%%%%%%%%

% nguồn: MaHoaNH_TuLuan_Lop10
\begin{bt}%[0D8H3-6]%[ma:4f4d6a3d]
	Chứng minh  đẳng thức $\dfrac{1}{k+1} \mathrm{C}_n^k=\dfrac{1}{n+1} \mathrm{C}_{n+1}^{k+1}$.
	 \loigiai{
		Ta có  $\dfrac{1}{k+1} \mathrm{C}_n^k=\dfrac{n!}{(k+1) k!(n-k)!}=\dfrac{(n+1)!}{(n+1)(k+1)![n+1-(k+1)]!}=\dfrac{1}{n+1} \mathrm{C}_{n+1}^{k+1}$.
	} 
\end{bt}

% nguồn: MaHoaNH_TuLuan_Lop10
\begin{bt}%[0D8H3-6]%[ma:686dc57f]
	Chứng  minh đẳng thức $k \mathrm{C}_n^k=n \mathrm{C}_{n-1}^{k-1}$.
	\loigiai{ 
		Ta có  $k \mathrm{C}_n^k=k\cdot \dfrac{n!}{k!(n-k)!}=\dfrac{n!(n-1)!}{(k-1)![n-1-(k-1)]!}=n \mathrm{C}_{n-1}^{k-1}$.
	 }
\end{bt}

% nguồn: MaHoaNH_TuLuan_Lop10
\begin{bt}%[0D8V3-6]%[ma:2ee1bae8]
	 Cho khai triển $(1+2x)^n=a_0+a_1x+a_2x^2+\ldots+a_nx^n$ thỏa mãn $a_0+8 a_1=2 a_2+1$. Tìm giá trị của số nguyên dương $n$.
	\loigiai{ 
		Ta có $(1+2x)^n=\sum\limits_{k=0}^n2^k\mathrm{C}_n^kx^k$; $(k \in \mathbb{N})$. Suy ra $a_k=2^k \mathrm{C}_n^k$. \\ 
		Thay $a_0=2^0\mathrm{C}_n^0=1$, $a_1=2 \mathrm{C}_n^1$, $a_2=4 \mathrm{C}_n^2$ vào giả thiết ta  có
		\begin{eqnarray*}
			1+16\mathrm{C}_n^1=8\mathrm{C}_n^2+1 
			 &\Leftrightarrow& 2\mathrm{C}_n^1=\mathrm{C}_n^2
			 \Leftrightarrow 2\dfrac{n!}{(n-1)!}=\dfrac{n!}{(n-2)!2!} \\
			&\Leftrightarrow&  2n=\dfrac{n(n-1)}{2}
			\Leftrightarrow n^2-5 n=0 \Leftrightarrow  \hoac{&n=0 \\ &n=5.}
		\end{eqnarray*}
		Do $n$ là số nguyên  dương nên $n=5$.
	} 
\end{bt}
